\begin{lemma}
\label{lemma-colimit-projective-dimension}
Let $R$ be a ring. Suppose we have a module $M = \bigcup_{e \in E} M_e$
where the $M_e$ are submodules well-ordered by inclusion. Assume the quotients
$M_e/\bigcup\nolimits_{e' < e} M_{e'}$ have projective dimension $\leq n$.
Then $M$ has projective dimension $\leq n$.
\end{lemma}

\begin{proof}
We will prove this by induction on $n$.

\medskip\noindent
Base case: $n = 0$. Then $P_e = M_e/\bigcup_{e' < e} M_{e'}$ is projective.
Thus we may choose a section $P_e \to M_e$ of the projection $M_e \to P_e$.
We claim that the induced map $\psi : \bigoplus_{e \in E} P_e \to M$ is an
isomorphism. Namely, if $x = \sum x_e \in \bigoplus P_e$ is nonzero,
then we let $e_{max}$ be maximal such that $x_{e_{max}}$ is nonzero
and we conclude that $y = \psi(x) = \psi(\sum x_e)$ is nonzero because
$y \in M_{e_{max}}$ has nonzero image $x_{e_{max}}$ in $P_{e_{max}}$.
On the other hand, let $y \in M$. Then $y \in M_e$ for some $e$.
We show that $y \in \Im(\psi)$ by transfinite induction on $e$.
Let $x_e \in P_e$ be the image of $y$. Then
$y - \psi(x_e) \in \bigcup_{e' < e} M_{e'}$.
By induction hypothesis we conclude that $y - \psi(x_e) \in \Im(\psi)$
hence $y \in \Im(\psi)$. Thus the claim is true and
$\psi$ is an isomorphism. We conclude that $M$ is projective as
a direct sum of projectives, see
Lemma \ref{lemma-direct-sum-projective}.

\medskip\noindent
If $n > 0$, then for $e \in E$ we denote $F_e$ the free $R$-module
on the set of elements of $M_e$. Then we have a system of
short exact sequences
$$
0 \to K_e \to F_e \to M_e \to 0
$$
over the well-ordered set $E$. Note that the transition maps
$F_{e'} \to F_e$ and $K_{e'} \to K_e$ are injective too.
Set $F = \bigcup F_e$ and $K = \bigcup K_e$. Then
$$
0 \to
K_e/\bigcup\nolimits_{e' < e} K_{e'} \to
F_e/\bigcup\nolimits_{e' < e} F_{e'} \to
M_e/\bigcup\nolimits_{e' < e} M_{e'} \to 0
$$
is a short exact sequence of $R$-modules too and
$F_e/\bigcup_{e' < e} F_{e'}$ is the free $R$-module on the
set of elements in $M_e$ which are not contained in $\bigcup_{e' < e} M_{e'}$.
Hence by
Lemma \ref{lemma-exact-sequence-projective-dimension}
we see that the projective dimension of $K_e/\bigcup_{e' < e} K_{e'}$
is at most $n - 1$. By induction we conclude that $K$ has projective
dimension at most $n - 1$. Whence $M$ has projective dimension at most
$n$ and we win.
\end{proof}